\documentclass[11pt]{article}
\usepackage[T1]{fontenc}
\usepackage{lmodern}
\usepackage{microtype}
\usepackage[a4paper,margin=25mm]{geometry}
\usepackage{amsmath,amssymb,amsthm,mathtools}
\usepackage{enumitem}
\usepackage[hidelinks]{hyperref}
\usepackage{fancyhdr}
\hypersetup{pdftitle={Smooth-gradient approximation of finite-Fisher-information scores},pdfsubject={A proof of AIM catalogue Problem 618},pdfauthor={Research note prepared in this conversation}}
\setlength{\parskip}{3pt}
\setlength{\parindent}{0pt}
\setlength{\headheight}{14pt}
\pagestyle{fancy}
\fancyhf{}
\fancyhead[L]{\small Finite-Fisher-information scores}
\fancyhead[R]{\small AIM Problem 618}
\fancyfoot[C]{\thepage}
\renewcommand{\headrulewidth}{0.3pt}
\newtheorem{theorem}{Theorem}
\newtheorem{corollary}[theorem]{Corollary}
\newcommand{\dd}{\,\mathrm{d}}
\newcommand{\dV}{\dd V_g}
\newcommand{\I}{\mathcal I}
\newcommand{\normrho}[1]{\left\lVert #1\right\rVert_{L^2(\rho\,\mathrm{d}V_g)}}
\newcommand{\iprho}[2]{\left\langle #1,#2\right\rangle_{L^2(\rho\,\mathrm{d}V_g)}}
\title{\vspace{-10mm}\textbf{Smooth-gradient approximation of\\ finite-Fisher-information scores}\\[5pt]\large A proof of AIM catalogue Problem 618}
\author{}
\date{2 October 2026}
\begin{document}
\thispagestyle{plain}
\begin{center}
{\Large\bfseries Smooth-gradient approximation of\\[2pt]
finite-Fisher-information scores}\par
\vspace{6pt}
{\large A proof of AIM catalogue Problem 618}\par
\vspace{10pt}
2 October 2026
\end{center}
\vspace{2pt}
\begin{abstract}
Let $M$ be a smooth compact Riemannian manifold, possibly with smooth boundary, and let $\rho$ be a nonnegative probability density with $\sqrt\rho\in H^1(M)$. We prove that its score is in the $L^2(\rho\,\mathrm{d}V_g)$ closure of smooth gradients. An explicit approximation is
\[
 \varphi_\tau=P_\tau\log(P_\tau\rho+\tau),\qquad \tau\downarrow0,
\]
where $P_\tau$ is the Neumann heat semigroup. These potentials are smooth up to the boundary and have zero normal derivative. The proof uses a short-time heat-semigroup gradient estimate, an exact mixed-product identity, and lower semicontinuity of the square-root Dirichlet energy. It does not assume strict positivity or boundedness of $\rho$, or convexity of the boundary.
\end{abstract}

\section{Statement and analytic input}
The target is Problem 618 of the AIM catalogue \cite{AIM}, corresponding to the conjecture following Theorem 4 of Casteras--Flaim--Monsaingeon \cite{CFM}.

\begin{theorem}\label{thm:main}
Let $(M,g)$ be a smooth compact Riemannian manifold, possibly disconnected and possibly with smooth boundary. Suppose
\[
 \rho\geq0,\qquad \int_M\rho\dV=1,\qquad u:=\sqrt\rho\in H^1(M).
\]
Set
\[
 s_\rho=\begin{cases}2\nabla u/u,&u>0,\\0,&u=0,\end{cases}
 \qquad \I(\rho)=4\int_M|\nabla u|_g^2\dV.
\]
Let $P_\tau=\exp(\tau\Delta_N)$ be the heat semigroup with homogeneous Neumann boundary condition; on a boundaryless manifold use the ordinary heat semigroup. Then
\begin{equation}\label{eq:construction}
 \varphi_\tau=P_\tau\!\left[\log\big(P_\tau\rho+\tau\big)\right]
 \in C^\infty(M),\qquad \tau>0,
\end{equation}
where smoothness is understood up to the boundary, and
\begin{equation}\label{eq:conclusion}
 \lim_{\tau\downarrow0}\int_M|\nabla\varphi_\tau-s_\rho|_g^2\rho\dV=0.
\end{equation}
Moreover, $\partial_\nu\varphi_\tau=0$ on $\partial M$. In particular, $\varphi_j:=\varphi_{1/j}$ supplies the sequence requested by the catalogue problem.
\end{theorem}

We write $\Delta=\operatorname{div}\nabla$, so $\Delta_N$ is nonpositive on $L^2(M)$. The nontrivial analytic input is the following established estimate: there are $\tau_0>0$ and a function $c:(0,\tau_0]\to[1,\infty)$, independent of $f$, such that
\begin{equation}\label{eq:GE}
 |\nabla P_\tau f|_g^2\leq c(\tau)P_\tau(|\nabla f|_g^2),
 \qquad f\in C^1(M),\qquad \lim_{\tau\downarrow0}c(\tau)=1.
\end{equation}
For a smooth nonconvex boundary this is Sturm's Theorem 1.1(iii), with $p=2$, and Remark 1.2 \cite{Sturm}. If the second fundamental form is bounded below by $-S$, one can take, for small $\tau$,
\[
 c(\tau)=\exp\!\left(4S\sqrt{\tau/\pi}+C\tau\right)
\]
with $C\geq0$ depending on the geometry. On a closed manifold the usual Ricci-lower-bound estimate gives $c(\tau)=e^{2K\tau}$. The requisite geometric lower bounds exist by smoothness and compactness. No convexity is required. Estimate \eqref{eq:GE} is also recorded as (GE) in \cite[Section 3]{CFM}.

We use the standard properties of this semigroup: positivity, preservation of constants, self-adjointness with respect to $\mathrm{d}V_g$, strong continuity on $L^1(M)$, and smoothing up to the boundary at positive times. Also
\begin{equation}\label{eq:commute}
 \Delta_N P_\tau h=P_\tau\Delta_N h\qquad (h\in D(\Delta_N)).
\end{equation}
Importantly, the functions $h$ to which \eqref{eq:commute} is applied below are smooth and satisfy the Neumann condition.

\section{Proof}
Throughout the proof norms and scalar products of tangent vectors are taken with respect to $g$.

\subsection{Sobolev identities for the original density}
The product rule for Sobolev functions implies
\begin{equation}\label{eq:rhoW11}
 \rho=u^2\in W^{1,1}(M),\qquad \nabla\rho=2u\nabla u\in L^1(M;TM).
\end{equation}
For completeness, if $u_n\to u$ in $H^1(M)$ with $u_n$ smooth, then
\[
 \|u_n^2-u^2\|_{L^1}
 \leq(\|u_n\|_{L^2}+\|u\|_{L^2})\|u_n-u\|_{L^2}\longrightarrow0,
\]
and
\begin{align*}
 \|u_n\nabla u_n-u\nabla u\|_{L^1}
 &\leq\|u_n-u\|_{L^2}\|\nabla u_n\|_{L^2}
       +\|u\|_{L^2}\|\nabla u_n-\nabla u\|_{L^2}\\
 &\longrightarrow0.
\end{align*}
This proves \eqref{eq:rhoW11} distributionally. The locality property of Sobolev gradients gives $\nabla u=0$ almost everywhere on $\{u=0\}$. Consequently,
\begin{equation}\label{eq:scoreidentity}
 \rho s_\rho=\nabla\rho,\qquad
 \normrho{s_\rho}^{\,2}=4\int_M|\nabla u|^2\dV=\I(\rho).
\end{equation}
No assertion that $\log\rho$ belongs to an unweighted Sobolev space is being made or used.

For every smooth $\psi$ satisfying $\partial_\nu\psi=0$, Green's identity gives
\begin{equation}\label{eq:weakgreen}
 \int_M\rho s_\rho\cdot\nabla\psi\dV
 =\int_M\nabla\rho\cdot\nabla\psi\dV
 =-\int_M\rho\Delta\psi\dV.
\end{equation}
This remains valid for $\rho\in W^{1,1}$: it follows either from the trace form of Green's identity or by using $u_n^2\to\rho$ in $W^{1,1}$ as above. No boundary condition on $\rho$ is needed.

\subsection{Regularization and the exact mixed product}
Fix $\tau>0$ and set
\[
 r_\tau=P_\tau\rho,\qquad q_\tau=r_\tau+\tau,
 \qquad h_\tau=\log q_\tau,\qquad \varphi_\tau=P_\tau h_\tau.
\]
The function $r_\tau$ is smooth and nonnegative, and $q_\tau\geq\tau>0$. Thus $h_\tau$ and $\varphi_\tau$ are smooth up to the boundary. Furthermore,
\[
 \partial_\nu r_\tau=0,
 \qquad \partial_\nu h_\tau=
 \frac{\partial_\nu r_\tau}{r_\tau+\tau}=0,
 \qquad \partial_\nu\varphi_\tau=0.
\]
In particular $h_\tau\in D(\Delta_N)$, so \eqref{eq:commute} applies.

Define two finite nonnegative quantities
\begin{equation}\label{eq:AB}
 A_\tau=\int_M\frac{|\nabla r_\tau|^2}{r_\tau+\tau}\dV,
 \qquad B_\tau=\int_M\rho|\nabla\varphi_\tau|^2\dV.
\end{equation}
The key identity is
\begin{equation}\label{eq:mixed}
 \iprho{s_\rho}{\nabla\varphi_\tau}=A_\tau.
\end{equation}
Indeed, using \eqref{eq:weakgreen}, commutation, and self-adjointness successively,
\begin{align*}
 \int_M\rho s_\rho\cdot\nabla\varphi_\tau\dV
 &=\int_M\nabla\rho\cdot\nabla P_\tau h_\tau\dV\\
 &=-\int_M\rho\Delta P_\tau h_\tau\dV\\
 &=-\int_M\rho P_\tau\Delta h_\tau\dV\\
 &=-\int_M r_\tau\Delta h_\tau\dV\\
 &=\int_M\nabla r_\tau\cdot\nabla h_\tau\dV\\
 &=\int_M\frac{|\nabla r_\tau|^2}{r_\tau+\tau}\dV=A_\tau.
\end{align*}
The first integration by parts uses $\partial_\nu\varphi_\tau=0$; the second uses $\partial_\nu h_\tau=0$. Self-adjointness here is valid in the $L^1$--$L^\infty$ pairing, since $\rho\in L^1$ and $\Delta h_\tau$ is bounded for each fixed $\tau$.

\subsection{The fixed-weight norm bound}
Applying \eqref{eq:GE} to $h_\tau$, multiplying by the \emph{original} density $\rho$, and integrating yields
\begin{align}
 B_\tau
 &\leq c(\tau)\int_M\rho P_\tau(|\nabla h_\tau|^2)\dV\notag\\
 &=c(\tau)\int_M r_\tau|\nabla h_\tau|^2\dV\notag\\
 &=c(\tau)\int_M\frac{r_\tau|\nabla r_\tau|^2}{(r_\tau+\tau)^2}\dV\notag\\
 &\leq c(\tau)A_\tau.\label{eq:Bbound}
\end{align}
Uniformity of $c(\tau)$ in the input function is essential, since $h_\tau$ depends on $\tau$.

Write $I=\I(\rho)$. By \eqref{eq:scoreidentity}, \eqref{eq:mixed}, and Cauchy--Schwarz in $L^2(\rho\,\mathrm{d}V_g)$,
\[
 A_\tau^2\leq I B_\tau\leq I c(\tau)A_\tau.
\]
If $A_\tau>0$, division by $A_\tau$ gives $A_\tau\leq c(\tau)I$. The same conclusion is immediate if $A_\tau=0$. We have proved
\begin{equation}\label{eq:Abound}
 0\leq A_\tau\leq c(\tau)I.
\end{equation}
This bound was obtained directly from the original finite Fisher information; it does not presuppose the approximation theorem.

\subsection{Convergence of the regularized square-root energies}
Set $v_\tau=\sqrt{r_\tau+\tau}$. Strong continuity of $P_\tau$ on $L^1(M)$ gives
\[
 \|r_\tau-\rho\|_{L^1}\longrightarrow0.
\]
Since $|\sqrt a-\sqrt b|^2\leq|a-b|$ for $a,b\geq0$,
\begin{equation}\label{eq:L2}
 \|v_\tau-u\|_{L^2}^2
 \leq\|r_\tau+\tau-\rho\|_{L^1}
 \leq\|r_\tau-\rho\|_{L^1}+\tau\operatorname{Vol}_g(M)
 \longrightarrow0.
\end{equation}
On the other hand,
\begin{equation}\label{eq:energy}
 4\int_M|\nabla v_\tau|^2\dV=A_\tau\leq c(\tau)I.
\end{equation}
Thus $v_\tau$ is bounded in $H^1(M)$ for small $\tau$. Weak lower semicontinuity of its Dirichlet energy, together with \eqref{eq:L2}, implies
\[
 I=4\int_M|\nabla u|^2\dV\leq\liminf_{\tau\downarrow0}A_\tau.
\]
Explicitly, apply weak compactness to any sequence of times realizing the lower limit; its weak $H^1$ limit is $u$ by the strong $L^2$ convergence. Conversely, \eqref{eq:Abound} and $c(\tau)\to1$ imply
\[
 \limsup_{\tau\downarrow0}A_\tau\leq I.
\]
Therefore
\begin{equation}\label{eq:Aconvergence}
 A_\tau\longrightarrow I.
\end{equation}

\subsection{Strong convergence in the required weighted space}
Expand the squared error using \eqref{eq:mixed}:
\begin{align*}
 \normrho{\nabla\varphi_\tau-s_\rho}^{\,2}
 &=B_\tau+I-2A_\tau\\
 &\leq I+\big(c(\tau)-2\big)A_\tau.
\end{align*}
By \eqref{eq:Aconvergence} and $c(\tau)\to1$, the right-hand side tends to zero. The left-hand side is nonnegative, so \eqref{eq:conclusion} follows. All asserted smoothness and boundary properties were established above. This proves Theorem~\ref{thm:main}. \qed

\section{Scope, mechanism, and a flux consequence}
\subsection{Why none of the excluded assumptions has been reintroduced}
The shift $+\tau$ is applied only to the regularization, not to the given density. It makes the logarithm well-defined even on a connected component carrying zero mass. It also permits arbitrary vacuum regions in the original density.

An upper bound on $\rho$ is never used. The only facts needed to handle the original density are $\rho\in W^{1,1}$, the identity $\rho s_\rho=\nabla\rho$, and $s_\rho\in L^2(\rho\,\mathrm{d}V_g)$. Every smooth function used in an integration by parts is fixed at a positive heat time.

Nonconvexity is handled solely through the factor $c(\tau)=1+O(\sqrt\tau)$. The proof requires that this factor tends to one, not that it is identically one or has an $O(\tau)$ error. The input density has no imposed boundary condition. The homogeneous Neumann condition belongs to the constructed test potentials and is allowed by the problem.

\subsection{Why the second application of the heat semigroup matters}
Regularizing $\rho$ alone controls expressions weighted by $P_\tau\rho$, whereas the requested convergence is weighted by $\rho$. The outer $P_\tau$ in \eqref{eq:construction} makes self-adjointness transfer the regularized weight back to the original one in \eqref{eq:Bbound}; it also gives the exact mixed product \eqref{eq:mixed}. No general weighted-mollification density assertion is used.

The double-regularization calculation is suggested by the proof of Theorem 2 in \cite{CFM}. The conclusion here follows by retaining the mixed product, combining it with Cauchy--Schwarz, and applying lower semicontinuity before expanding the fixed-weight error.

\begin{corollary}[The score annihilates invisible fluxes]\label{cor:flux}
Under the hypotheses of Theorem~\ref{thm:main}, let $z\in L^2(\rho\,\mathrm{d}V_g;TM)$ satisfy
\[
 \int_M\rho z\cdot\nabla\psi\dV=0
 \qquad\text{for every }\psi\in C^1(M).
\]
Then
\[
 \int_M\rho z\cdot s_\rho\dV=0.
\]
\end{corollary}
\begin{proof}
For the potentials from Theorem~\ref{thm:main},
\begin{align*}
 \left|\int_M\rho z\cdot s_\rho\dV\right|
 &=\left|\int_M\rho z\cdot(s_\rho-\nabla\varphi_j)\dV\right|\\
 &\leq\normrho{z}\,\normrho{s_\rho-\nabla\varphi_j}
 \longrightarrow0.
\end{align*}
\end{proof}
Thus two finite-action velocities defining the same weak flux divergence and zero normal flux have the same pairing with the score. In particular, the additional gradient-closure hypothesis in the second conclusion of \cite[Theorem 4]{CFM} is automatic under its finite-Fisher-information assumptions. This is a consequence for the smooth geometry in that theorem; no assertion about the separate rough Lipschitz-domain chain-rule problem is made here.

\paragraph{Verification status.}
This is a proof note prepared in this conversation, not an independently refereed publication or an assertion that the catalogue's status has been changed.

\begin{thebibliography}{9}
\bibitem{AIM}
MColbrook/AIM, \emph{Problem 618: Smooth-gradient approximation of finite-Fisher-information scores}, in the ``PDEs, fluids and materials'' category. Catalogue page last checked 24 September 2026; accessed 2 October 2026.\newline
\href{https://github.com/MColbrook/AIM/blob/main/problems/618-fisher-score-gradient-closure.md}{Repository problem statement}.

\bibitem{CFM}
J.-B. Casteras, M. Flaim, and L. Monsaingeon,
\emph{Entropy and Fisher information in non-convex domains: one chain to rule them all},
\href{https://arxiv.org/abs/2512.05826v2}{arXiv:2512.05826v2}, 2026.
Conjecture following Theorem 4; Section 3, estimate (GE) and proof of Theorem 2; Lemma A.3.

\bibitem{Sturm}
K.-T. Sturm,
\emph{Gradient and Transport Estimates for Heat Flow on Nonconvex Domains},
\href{https://arxiv.org/abs/2502.01915v1}{arXiv:2502.01915v1}, 2025.
Theorem 1.1(iii), Remark 1.2, and Remark 1.3.
\end{thebibliography}
\end{document}
